\label{chap:p3-83-47-godel_cohen}

\section{From constructible consistency to independence by forcing}

\begin{theorem}[Gödel, 1940: constructible half]
In usual metatheory,
\[
  \operatorname{Con}(\mathsf{ZF})
  \Longrightarrow
  \operatorname{Con}
  \bigl(\mathsf{ZF}+\mathsf{AC}+\mathsf{GCH}\bigr).
\]
More precisely, given a model \(M\models\mathsf{ZF}\), its internal
constructible universe \(L^M\) satisfies, in the corresponding internal sense,
\[
  \mathsf{ZF}+V=L+\mathsf{AC}+\mathsf{GCH}.
\]
\end{theorem}

\begin{proof}[Idea of the classical proof]
The constructible hierarchy is defined by recursion over the ordinals, admitting at each stage only subsets definable over earlier stages. The construction preserves the ordinals, produces an inner model, and provides a global definable well-order, from which choice and the generalized continuum hypothesis \citep{goedel1940} follow.
\end{proof}

\begin{theorem}[Cohen, 1963--1964: forcing half]
One has the metamathematical implications
\[
  \operatorname{Con}(\mathsf{ZFC})
  \Longrightarrow
  \operatorname{Con}(\mathsf{ZFC}+\neg\mathsf{CH}),
\]
and
\[
  \operatorname{Con}(\mathsf{ZF})
  \Longrightarrow
  \operatorname{Con}(\mathsf{ZF}+\neg\mathsf{AC}).
\]
Combining the two halves, \(\mathsf{CH}\) is independent of
\(\mathsf{ZFC}\) and \(\mathsf{AC}\) is independent of \(\mathsf{ZF}\),
always under the corresponding consistency hypothesis
\citep{cohen1963,cohen1964}.
\end{theorem}

For the first construction, one begins with a base model \(M\), a poset
\(\mathbb P\in M\), and an \(M\)-generic filter \(G\subseteq\mathbb P\):
\[
  G\cap D\ne\varnothing
  \qquad
  \text{for every dense }D\subseteq\mathbb P\text{ belonging to }M.
\]
The forcing theorem relates \(p\Vdash_M\varphi\) to the truth of
\(\varphi\) in \(M[G]\). To obtain
\(\mathsf{ZF}+\neg\mathsf{AC}\), one additionally takes a suitable symmetric
submodel of a generic extension; an ordinary extension of a model of
\(\mathsf{ZFC}\) preserves choice \citep{cohen1966}.

\begin{remark}[Conditions of validity]
These theorems prove relative consistency and independence. They do not prove the
absolute consistency of \(\mathsf{ZF}\) or \(\mathsf{ZFC}\), do not decide
\(\mathsf{CH}\) in a privileged universe, and do not convert every compatible filter
into a generic filter. Nor do they select the HMT branches or
identify a numerical output.
\end{remark}

\subsection{\(2\) Non-Equivalent Operations: Model Extension and Object Enrichment}

Forcing and HMT do not perform the same operation. The former compares a
universe \(M\) with an extension \(M[G]\); the latter preserves fibers that a
reader forgets. If \(\mathsf{Hist}_{\rm HMT}\) is the category of enriched
histories, the functor
\[
  U:\mathsf{Hist}_{\rm HMT}\longrightarrow\mathsf{Set},
  \qquad
  (x_m,\mu_m)_{m\ge0}\longmapsto(x_m)_{m\ge0},
\]
may identify distinct histories without changing the visible set. Cohen's
question is which subsets exist in the model; the genealogical question is
how many lifts, with memory and operators, a visible object possesses. They
are compatible axes, but neither replaces the other.

\begin{center}
\begin{tabularx}{.94\textwidth}{>{\bfseries}p{.22\textwidth}XX}
\toprule
Operation & What changes & What does not prove by itself\\
\midrule
\(M\mapsto M[G]\) & Set-theoretic universe, names and forced truth &
Uniqueness or genealogy of an HMT section.\\
\(\xi\mapsto U(\xi)\) & Amount of visible memory in the projection &
\(\mathsf{CH}\), \(\mathsf{AC}\), or genericity over \(M\).\\
\bottomrule
\end{tabularx}
\end{center}

\subsection{HMT power does not presuppose \texorpdfstring{\(\mathsf{CH}\)}{CH}}

In every model in which \(X\) is a Cantor space,
\[
  |\Clop(X)|=\aleph_0,\qquad
  |\Closed(X)|=\mathfrak c,\qquad
  |\mathcal P(X)|=2^{\mathfrak c},
  \qquad \mathfrak c=2^{\aleph_0}.
\]
\(\mathsf{CH}\) asserts only \(\mathfrak c=\aleph_1\). By Cantor's theorem,
\(\mathfrak c<2^{\mathfrak c}\) in each of the models
being compared. On passing from \(M\) to \(M[G]\), new branches, new
closed codes, and new subsets may appear, although the finite cylinders
retain their formation rule. The full power set is model-sensitive; the
clopen--closed--full distinction does not depend on deciding \(\mathsf{CH}\).

\section{Compatible prefixes poset and inverse limit of extensions}

We define
\[
  \mathbb P_{\TPK}=\bigsqcup_{n\ge0}\mathcal S_n
\]
and order by prolongation: \(q\le p\) if \(q\) contains more information
than \(p\). For each \(N\),
\[
  D_N=\{p\in\mathbb P_{\TPK}:\ell(p)\ge N\}.
\]
If every surviving prefix can be prolonged, each \(D_N\) is dense.

\begin{theorem}[Branches and depth filters]
The sections of \(\OmegaInf\) are in bijection with the compatible filters
of \(\mathbb P_{\TPK}\) that intersect all the \(D_N\).
\end{theorem}

\begin{proof}
The prefixes of a section are compatible, directed, and reach every
depth. Conversely, a filter that meets each \(D_N\) determines a
unique restriction at each level; directedness excludes incompatible
conditions and produces a section.
\end{proof}

Meeting all dense sets \(D_N\) expresses unbounded depth, not
genericity over a model. Being \(M\)-generic would require intersecting every
dense set \(D\in M\).

\begin{proposition}[Conditioned Cohen case]
If \(\OmegaInf\) is compact, metrizable, zero-dimensional, perfect, and without
isolated points, it is homeomorphic to the Cantor set; the forcing of nonempty cylinders
has a dense part equivalent to Cohen forcing,
\citep{cohen1963,cohen1964}.
\end{proposition}

A unique section, by contrast, produces an atomic or trivial forcing.
The word “forcing” does not by itself select a branch.

For a section \(x=(x_m)_{m\ge0}\), its prefixes form a filter \(G_x\)
meeting the depth-dense sets. Nevertheless,
\[
  \bigl(\forall N,\;G_x\cap D_N\ne\varnothing\bigr)
  \centernot\Longrightarrow
  \bigl(\forall D\in M\text{ dense},\;G_x\cap D\ne\varnothing\bigr).
\]
The first property expresses unbounded prolongation; the second,
genericity relative to a model. A compatible section produced by
nine TPK phases does not become a Cohen real by nomenclature alone.

\begin{remark}[Conditions of validity]
The following falsify an automatic identification with Cohen: a terminal cylinder or an
isolated point; a filter that meets every \(D_N\) but omits some
dense subset of \(M\); two compatible memories on the same visible branch; or the
absence of an explicit dense embedding. The shared name “branch” is not
a morphism.
\end{remark}

\section{Correspondence between genealogical sections and Borel classes of the limit space}

The cylinders generate
\[
 \Gamma_{\rm HMT}^{\rm Bor}
 =
 \sigma\bigl(\{[p]_n:p\in\mathcal S_n\}\bigr).
\]
The class is stable under complements, countable unions, and preimages of certified continuous maps.

If an infinite game on histories has a Borel winning condition,
Martin's Borel determinacy theorem guarantees a winning strategy
\citep{martin1975}. This is a classical application to HMT space;
it does not derive AD, MM, a value of \(2^{\aleph_0}\), or categoricity of
\(\mathcal P(\mathbb R)\).

For \(A,B\subseteq\OmegaInf\), the Wadge reduction
\[
 A\le_WB
\]
requires a continuous map \(f\) with \(x\in A\iff f(x)\in B\). If the
maps are restricted to certified TPK operators, a narrower preorder appears,
whose well-foundedness or semilinearity requires a proof of its own.

\begin{remark}[Conditions of validity]
A proof that uses Wadge well-foundedness to establish the
determinacy that would then justify that well-foundedness is circular. The
historical AD/MM branches are retained as architecture, not as HMT theorems.
\end{remark}

\section{Minimax optimization over finite windows of compatible prefixes}

Let \(P_N=\mathcal S_N\) be the finite set of candidates, \(T_N\) a finite set
of adversarial tests, and
\[
  A_N:P_N\times T_N\to[-1,1]
\]
a payment fixed before the result.

\begin{theorem}[Finite minimax]
\[
 \max_{\sigma\in\Prob(P_N)}
 \min_{\tau\in\Prob(T_N)}
 \mathbb E_{\sigma,\tau}A_N
 =
 \min_{\tau\in\Prob(T_N)}
 \max_{\sigma\in\Prob(P_N)}
 \mathbb E_{\sigma,\tau}A_N.
\]
\end{theorem}

\begin{proof}
It is the strong duality between the selector's linear program and that of the
adversary, in the form of von Neumann's minimax theorem,
\citep{vonneumann1928}.
\end{proof}

Equilibrium may be mixed and does not provide a unique pure branch.

\section{Projective system of truncations and construction of the limit section}

Let
\[
 X=\varprojlim P_N,
 \qquad
 Y=\varprojlim T_N
\]
compact, and let \(A:X\times Y\to\mathbb R\) be continuous. The spaces of
probability measures are compact and convex in the weak topology. Sion's
theorem \citep{sion1958} yields
\[
 \sup_{\mu\in\Prob(X)}\inf_{\nu\in\Prob(Y)}
 \iint A\,d\mu\,d\nu
 =
 \inf_{\nu\in\Prob(Y)}\sup_{\mu\in\Prob(X)}
 \iint A\,d\mu\,d\nu.
\]

\begin{corollary}
A global equilibrium selects a pure branch only if the selector measure
is a Dirac measure \(\delta_x\). The existence of a mixed equilibrium does not imply
concentration.
\end{corollary}

\begin{remark}[Conditions of validity]
Passage to the limit fails without compactness, continuity, or projective
consistency of the measures. If the payoff is defined using the target it is
to select, the minimax is target-fed and does not explain the route.
\end{remark}

\section{Projective and Matrix Extensions Provided by Operator Entanglers}

Forcing organizes possible prolongations; Borel theory classifies sets of
histories; minimax constructs robustness against falsifiers. None replaces
the structural law that produces a unique survivor. Their function is to
clarify the topology and decision structure of route space, not to introduce a
retrospective selection.

\begin{remark}[Separation of results and scope]
Gödel's and Cohen's theorems are \emph{applied classical theorems}. The
poset, the typed comparison between the depth filter and the generic filter,
the HMT Borel class, and the selector--falsifier game are
\emph{explicit constructions of the present development}. Martin, von Neumann, and Sion are applied
thereafter, with their classical hypotheses. The former PSC/HTT assertions concerning
CH, AD, and MM are not promoted.
\end{remark}
